\subsection{The relation of domination between local rings}

DEFINITION 1. Let A and B be two local rings. B is said to dominate A if A is a subring of B and \(\mathfrak{m}(\mathbf{A}) = \mathbf{A} \cap \mathfrak{m}(\mathbf{B})\) .

PROPOSITION 1. Let A and B be local rings such that A is a subring of B. The following conditions are equivalent:

\begin{enumerate}
\def\labelenumi{(\alph{enumi})}
\item
  \(\mathfrak{m}(\mathbf{A}) \subset \mathfrak{m}(\mathbf{B});\)
\item
  B dominates A;
\item
  the ideal \(\mathbf{B}\mathfrak{m}(\mathbf{A})\) generated by \(\mathfrak{m}(\mathbf{A})\) in \(\mathbf{B}\) does not contain 1.
\end{enumerate}

If \(\mathfrak{m}(\mathbf{A}) \subset \mathfrak{m}(\mathbf{B})\) , \(\mathfrak{m}(\mathbf{B}) \cap \mathbf{A}\) is an ideal of A which does not contain 1 and which contains the maximal ideal \(\mathfrak{m}(\mathbf{A})\) ; it is therefore equal to it and (a) implies (b). If B dominates A, the ideal \(\mathbf{B}\mathfrak{m}(\mathbf{A})\) is contained in \(\mathfrak{m}(\mathbf{B})\) and hence does not contain 1; thus (b) implies (c). If (c) holds, \(\mathbf{B}\mathfrak{m}(\mathbf{A})\) is contained in the unique maximal ideal \(\mathfrak{m}(\mathbf{B})\) of B, whence (a).

Note that, if K is a ring, the relation ``B dominates A'' is an order relation on the set of local subrings of K.

Let A and B be two local rings such that B dominates A. The canonical injection A → B defines by taking quotients an isomorphism of the field κ(A) onto a subfield of κ(B); this isomorphism allows us to identify κ(A) with a subfield of κ(B).

\textbf{Examples}\par

\begin{enumerate}
\def\labelenumi{(\arabic{enumi})}
\item
  Let A be a Noetherian local ring and \(\hat{\mathbf{A}}\) its completion; the local ring \(\hat{\mathbf{A}}\) then dominates A (Chapter III, § 3, no. 5, Proposition 9).
\item
  Let B be an integral domain, A a subring of B, \(p'\) a prime ideal of B and \(p = A \cap p'\) . Then \(pA_{p} \subset p'B_{p'}\) , so that \(B_{p'}\) dominates \(A_{p}\) .
\end{enumerate}

